\chapter{Introduction}

\textbf{mpeg2fpga} is an MPEG-2 video decoder (ISO/IEC 13818-2) implemented in
the FPGA of a PolarFire SoC and offered as a network service. A client sends it
a clip over Ethernet and receives the decoded frames as they come out, in
planar YUV 4:2:0 (I420), in display order.

Decoding happens entirely in hardware. The SoC's RISC-V processors only
manage: they receive the clip, hand it to the decoder over DMA, read each
finished frame and send it to the client.

\section{What a release includes}

\begin{dato}
\begin{tabularx}{\linewidth}{@{}l X@{}}
\textbf{Device} & PolarFire SoC Discovery Kit board with the decoder
  bitstream programmed. The FPGA is flash-based: it keeps its configuration
  without power or external memory. \\[3pt]
\textbf{Board software} & Linux driver (\cmd{mpeg2fpga.ko}), the
  \cmd{mpeg2fpgad} service and its native frame converter. \\[3pt]
\textbf{Client library} & Python package \cmd{mpeg2fpga} (wheel), with no
  dependencies outside the standard library. \\[3pt]
\textbf{Protocol} & Specification of network protocol v1, the contract
  between the device and any client. \\[3pt]
\textbf{Documentation} & This guide and the release notes, with the
  conformance and performance report.
\end{tabularx}
\end{dato}

\section{Architecture}

\begin{center}
\begin{tikzpicture}[font=\sffamily\small,
    box/.style={blk, minimum width=31mm, minimum height=12mm},
    hbox/.style={hw, minimum width=31mm, minimum height=12mm}]
  % top row: software
  \node[box] (app) at (0,0) {\textbf{Application}\\[-1pt]\scriptsize \texttt{mpeg2fpga} library};
  \node[box] (d)   at (6.6,0) {\textbf{mpeg2fpgad}\\[-1pt]\scriptsize service, protocol v1};
  \node[box] (drv) at (11.2,0) {\textbf{Driver}\\[-1pt]\scriptsize \texttt{mpeg2fpga.ko}};
  % bottom row: FPGA
  \node[hbox] (fs)  at (2.0,-3.4) {\textbf{Frame store}\\[-1pt]\scriptsize DDR, 4 buffers};
  \node[hbox] (dec) at (6.6,-3.4) {\textbf{Decoder}\\[-1pt]\scriptsize MPEG-2};
  \node[hbox] (dma) at (11.2,-3.4) {\textbf{Input DMA}\\[-1pt]\scriptsize \texttt{stream\_dma}};
  % data path
  \draw[arr] ([yshift=2.5mm]app.east) -- node[above, lbl]{clip (HTTP)} ([yshift=2.5mm]d.west);
  \draw[arr] ([yshift=-2.5mm]d.west) -- node[below, lbl]{I420 frames} ([yshift=-2.5mm]app.east);
  \draw[arr] (d) -- node[above, lbl]{sysfs} (drv);
  \draw[arr] (drv) -- node[right, lbl]{clip chunks} (dma);
  \draw[arr] (dma) -- (dec);
  \draw[arr] (dec) -- node[above, lbl]{frames} (fs);
  % control path
  \draw[arr, dashed, color=m2teal] (dec.north east) to[bend right=12]
    node[pos=0.55, right=1mm, lbl, align=left]{IRQ:\\frame ready} (drv.south west);
  \draw[arr, dashed, color=m2teal] (fs.north) |- node[pos=0.3, left, lbl]{read} ([yshift=-6mm]d.south west) -- ([xshift=-8mm]d.south);
  \begin{scope}[on background layer]
    \node[draw=m2grey!60, dashed, rounded corners=4pt, fit=(d)(drv)(fs)(dma), inner sep=4mm] (soc) {};
    \node[lbl, anchor=south east] at (soc.north east) {PolarFire SoC};
    \node[lbl, anchor=south west] at ([yshift=6mm]app.north west) {Client PC};
  \end{scope}
\end{tikzpicture}
\end{center}

\begin{itemize}
  \item The \textbf{clip goes up} over HTTP. The service splits it into
    chunks and the driver hands them to the DMA, which feeds the decoder.
  \item When the decoder finishes a frame, an \textbf{interrupt} tells the
    driver which frame store buffer holds it, and the driver passes that on
    to the service. The service reads it at once, before the decoder reuses
    the buffer.
  \item The frame is converted to I420 and \textbf{comes down} to the client
    on the same connection, while the rest of the clip is still going up.
\end{itemize}

Upload and download happen at the same time, with flow control in both
directions: there is no limit on clip length and the memory used on the board
is fixed.

\section{Conventions in this guide}

Commands run on the user's PC have no prefix; commands run on the board are
prefixed with \cmd{ssh root@board}. In the examples the board's address is
\cmd{192.168.18.5}.
